% Owner: Kashaf Ali / Member A. Complete rewrite authorized 5 October 2026.
% Local font encoding restores the template's Times regular/bold faces under XeTeX.
\begingroup
\renewcommand{\encodingdefault}{T1}
\fontencoding{T1}\selectfont
\chapter{Project Vision}

Zeest brings patient history, consultation documentation and clinical analysis into a physician-controlled workspace. The physician selects one of their patients and an encounter, reviews relevant history and records the current consultation. Structured notes and supported imaging, laboratory and medication inputs provide the case information. Specialised processing and retrieved guidelines support a diagnosis and treatment draft, which is reviewed by the critic agent before presentation. The physician then decides whether to approve, modify or reject it. The vision is to preserve continuity across encounters while making findings, evidence and decision states understandable. Its value depends on the integration of existing capabilities and on examining their contributions through controlled comparisons.

\section{Problem Domain Overview}

Clinical decision support assists physicians in interpreting patient information and considering possible diagnoses or treatment plans. Zeest operates within this domain by connecting the current consultation with longitudinal history and supported clinical inputs. The Patient Memory Agent retrieves relevant context, while the Clinical Case Agent organises it with transcribed notes and other findings. Imaging, Laboratory and Medication Safety agents analyse their respective inputs. Retrieved guidelines support clinical reasoning, followed by critic review and physician inspection. The application provides registration, login and case navigation around this workflow. These functions support access and review; responsibility for accepting a recommendation remains with the physician rather than the processing agents.

\section{Problem Statement}

Patient information is distributed across encounters and represented through notes, images, laboratory findings and medication records. An isolated input does not establish its relationship to earlier observations or to the current consultation. Manual documentation adds preparation work, while generated recommendations are difficult to assess without their supporting evidence and clinical context. Zeest must connect historical retrieval, consultation capture, specialised analysis and evidence-supported reasoning without confusing a generated draft with an accepted decision. The problem therefore concerns both information continuity and accountable review. A coordinated workflow must preserve each input's source and encounter, present inspectable recommendations and require physician approval before recording a clinical decision.

\section{Problem Elaboration}

Longitudinal continuity requires relevant retrieval rather than simply storing every record. Historical observations must retain their patient, encounter and source relationships so that they can inform the correct case. Relational storage preserves these explicit associations, while semantic retrieval identifies context related to the current information need. Multimodal inputs introduce another challenge: consultation audio needs transcription, images need model-specific processing, and laboratory and medication information need suitable representations. The Clinical Case Agent connects these sources, and specialist agents contribute identifiable findings. Missing or unusable inputs must remain distinguishable from successful analysis so that incomplete processing is not mistaken for an established clinical result.

A clinical draft must also remain distinguishable from the evidence supporting it and the physician decision that follows. Retrieved guideline passages provide inspectable support, while the Safety and Critic Agent reviews the generated content before presentation. Modification changes a draft without approving it, rejection excludes it from accepted decisions, and approval records the particular revision reviewed by the physician. Audit history preserves these actions and their encounter context. The integrated workflow makes it possible to investigate whether memory, agents, critic review and multimodal inputs contribute useful support. Such comparisons require consistent cases and assessment conditions rather than assuming that additional components necessarily improve recommendations.

\section{Goals and Objectives}

The goal is to develop a physician-controlled medical copilot with longitudinal patient memory and coordinated multimodal reasoning. Consultation transcription supplies structured notes, specialist agents analyse supported inputs, and evidence retrieval grounds the generated diagnosis and treatment draft. Critic review and explicit physician decisions complete the workflow. Table~\ref{tab:kashaf:objectives} lists the nine objectives covering this integration and its comparative evaluation. Memory and transcription establish the information foundation; analysis, retrieval and reasoning prepare the recommendation; and approval controls preserve physician authority. The evaluation compares the single-language-model baseline, memory-augmented variant and full Zeest system to assess their contributions under consistent, clearly defined clinical assessment conditions.

\begin{table}[htbp]
\centering
\small
\caption{Zeest objectives covering memory, clinical processing, physician review and comparative evaluation.}
\label{tab:kashaf:objectives}
\begin{tabular}{|p{0.08\textwidth}|p{0.82\textwidth}|}
\hline
ID & Objective \\ \hline
O1 & Develop hybrid longitudinal memory using relational storage and semantic retrieval. \\ \hline
O2 & Convert consultation audio into structured clinical notes in real time. \\ \hline
O3 & Integrate Patient Memory, Clinical Case, Imaging, Laboratory and Medication Safety agents. \\ \hline
O4 & Integrate an existing validated pre-trained imaging classification model. \\ \hline
O5 & Retrieve relevant clinical guidelines and present supporting citations. \\ \hline
O6 & Generate a draft diagnosis for physician review within the clinical case workflow. \\ \hline
O7 & Review generated recommendations through the Safety and Critic Agent before presentation. \\ \hline
O8 & Provide approval, modification and rejection controls with a full audit history. \\ \hline
O9 & Compare a single LLM, a memory-augmented variant and the full Zeest system. \\ \hline
\end{tabular}
\end{table}

% Finish this section's floats before the next heading.
\clearpage
\section{Project Scope}

Zeest includes physician registration and access, patient navigation, encounter review, consultation transcription and structured-note correction. Each physician accesses only their own patient records. Its clinical workflow combines hybrid memory, specialised agents, supported imaging and laboratory inputs, medication context, evidence retrieval and critic review. Generated diagnoses and treatment plans require physician approval before recording. Development and evaluation use de-identified or synthetic cases, and clinical domain coverage depends on suitable data and existing validated models. Table~\ref{tab:kashaf:scope} summarises the boundaries. Training diagnostic models from scratch, autonomous treatment, emergency and ICU systems, real-world deployment, regulatory approval, native mobile applications, patient-facing chat and the optional knowledge graph are excluded from the project scope.

\begin{table}[htbp]
\centering
\small
\caption{Included capabilities and exclusions defining the boundaries of Zeest.}
\label{tab:kashaf:scope}
\begin{tabular}{|p{0.43\textwidth}|p{0.47\textwidth}|}
\hline
Included & Excluded \\ \hline
Physician registration, access and patient/encounter navigation & Separate administrator application \\ \hline
Consultation transcription and note review & Patient-facing chatbot and native mobile applications \\ \hline
Hybrid memory and specialised agents & Knowledge graph extension \\ \hline
Existing imaging models and clinical knowledge resources & Training diagnostic models from scratch \\ \hline
Evidence-supported drafts, critic review and physician decisions & Autonomous prescribing or treatment \\ \hline
De-identified/synthetic cases and supported data/model domains & Unsupported domains, emergency/ICU systems and real-world deployment \\ \hline
Controlled comparison of system variants & Regulatory approval and guaranteed clinical outcomes \\ \hline
\end{tabular}
\end{table}

% Finish this section's floats before the next heading.
\clearpage
\section{Sustainable Development Goal (SDG)}

Zeest aligns with Sustainable Development Goal 3, Good Health and Well-being, which promotes healthy lives and well-being for all ages.\cite{zeestA:unSDG3} Its intended contribution is to support physicians through clearer access to patient history, structured consultation documentation and reviewable clinical recommendations. Longitudinal memory connects earlier observations with the current encounter, while supporting evidence and critic findings help the physician inspect the generated draft. Explicit approval keeps the clinical decision under professional judgement. Figure~\ref{fig:1} shows Goal 3 within the seventeen Sustainable Development Goals. This alignment guides the focus on physician assistance; health-outcome improvements remain a matter for evaluation rather than an assumed result.

\begin{figure}[htbp]
\centering
\includegraphics[width=0.8\textwidth]{Figures/Picture1}
\caption{The seventeen Sustainable Development Goals shown as numbered coloured tiles; Goal 3, Good Health and Well-being, is the selected alignment for Zeest.}
\label{fig:1}
\end{figure}

% Finish this section's floats before the next heading.
\clearpage
\section{Constraints}

Clinical coverage depends on suitable case data, validated existing models and permitted resource use. Imaging integration must respect the selected model's input requirements and supported labels. Transcription depends on language coverage and audio conditions, while guideline retrieval and medication review require appropriate knowledge resources. Consistent interfaces must connect the backend, orchestration frameworks, language model and memory layer. The division between LangGraph and Google ADK requires detailed design agreement, and the embedding model must be compatible with the semantic store. Hardware and service arrangements affect performance, so numerical limits need a defined environment. Physician approval remains mandatory regardless of the selected components or supported clinical domain.

\section{Business Opportunity}

Zeest offers a coordinated workspace for consultation documentation, case preparation and recommendation review. Historical information, specialist findings and supporting guidelines can be considered within the same encounter instead of as disconnected outputs. The intended benefit is clearer access to clinical context and a more understandable review process. Integrating existing models and knowledge resources also provides a foundation that can accommodate compatible capabilities without rebuilding the entire workflow. For researchers and developers, the opportunity lies in identifying which components provide useful support and which add complexity. Future adoption depends on evidence of that value and the requirements of the intended setting, rather than an assumed commercial deployment or pricing model.

\section{Stakeholders Description/ User Characteristics}

Physicians are the direct users and reviewers of Zeest. They need readable notes, clear patient context and an understandable relationship between findings, evidence and generated recommendations. Clinical knowledge supports their decisions, while basic familiarity with a web interface supports navigation and input tasks. Patients are affected stakeholders whose interests concern appropriate handling of their information and physician oversight of recorded decisions. The student team develops the integration and maintains consistent component responsibilities. The supervisor and evaluation panel guide and assess the problem, design and research approach. These groups share an interest in a coherent, accountable workflow while retaining distinct clinical, development and academic responsibilities.

\subsection{Stakeholders Summary}

Table~\ref{tab:kashaf:stakeholders} summarises the responsibilities and interests of the principal stakeholders.

\begin{table}[htbp]
\centering
\small
\caption{Stakeholder responsibilities and principal interests in Zeest.}
\label{tab:kashaf:stakeholders}
\begin{tabular}{|p{0.2\textwidth}|p{0.28\textwidth}|p{0.39\textwidth}|}
\hline
Stakeholder & Responsibility/relationship & Principal interest \\ \hline
Physician & Direct user and reviewer of clinical recommendations & Relevant context, clear findings, inspectable evidence and control over recorded decisions. \\ \hline
Patient & Subject of clinical information and affected stakeholder & Appropriate information handling and physician oversight of recommendations. \\ \hline
Student team & System development and research integration & Consistent requirements, component interfaces and reproducible comparisons. \\ \hline
Supervisor and evaluation panel & Academic guidance and assessment & Justified scope, technical decisions and evidence addressing the research questions. \\ \hline
\end{tabular}
\end{table}

% Finish this section's floats before the next heading.
\clearpage
\subsection{Key High-Level Goals and Problems of Stakeholders}

Physicians need relevant history, reliable consultation notes and inspectable findings without losing the current patient context. Their main challenge is to assess information from several sources and determine whether it supports the generated draft. Patients need their information associated with the correct case and decisions reviewed by the responsible physician. The development team must connect processing components without creating ambiguous outputs or approval states. Academic reviewers need controlled comparisons that make the contributions of memory, agents, critic review and multimodal information assessable. Zeest addresses these needs through case continuity, source relationships and explicit review controls, with evaluation determining the usefulness of the integrated approach.

\endgroup
